\paragraph{折叠同余半环的合同格、端同态半格、幂零厚度与四相算术几何}
沿用附录定理 \ref{thm:fold-congruence-mod-semirings} 的模半环表述，记
\[
M_m:=F_{m+2},
\qquad
\mathcal R_m:=\Omega_m/\!\equiv_m,
\qquad
N(\omega):=\sum_{k=1}^{m}\omega_kF_{k+1},
\qquad
\pi_m(\omega):=N(\omega)\bmod M_m.
\]
又记
\[
\bar\pi_m:\mathcal R_m\xrightarrow{\ \sim\ }\ZZ/(M_m\ZZ)
\]
为由 \(\pi_m\) 诱导的半环同构。则折叠同余层上的全部合同、非幺端同态、幂零方向、局部/约化相图与 \(p\)-进可见深度，都可在单个 Fibonacci 模数 \(M_m\) 的约数理论与素因子型中完成而无信息损失。

\begin{theorem}[折叠同余半环的全部合同格与 Fibonacci 约数格]\label{thm:xi-fold-congruence-quotient-divisor-lattice}
\leanverified{paper\_xi\_fold\_congruence\_quotient\_divisor\_lattice}
对任意约数 \(d\mid M_m\)，令
\[
\rho_d:\ZZ/(M_m\ZZ)\longrightarrow \ZZ/(d\ZZ),
\qquad
\overline n\longmapsto \overline n \pmod d,
\]
并在 \(\mathcal R_m\) 上定义关系
\[
x\,\Theta_d\,y
\quad\Longleftrightarrow\quad
\rho_d\!\bigl(\bar\pi_m(x)\bigr)=\rho_d\!\bigl(\bar\pi_m(y)\bigr).
\]
则：
\begin{enumerate}
  \item \(\Theta_d\) 是 \(\mathcal R_m\) 的半环合同；
  \item 商半环满足规范同构
  \[
  \mathcal R_m/\Theta_d\ \cong\ \ZZ/(d\ZZ);
  \]
  \item \(\mathcal R_m\) 的任意半环合同都唯一地以某个 \(\Theta_d\) 的形式出现；
  \item 在合同格偏序下，有
  \[
  \Theta_{d_1}\le \Theta_{d_2}
  \iff
  d_2\mid d_1.
  \]
\end{enumerate}
因此，\(\operatorname{Con}(\mathcal R_m)\) 与 \(M_m\) 的约数格在整除偏序下反同构。
\end{theorem}
\begin{proof}
由于 \(d\mid M_m\)，若 \(n\equiv n'\pmod{M_m}\)，则必有 \(n\equiv n'\pmod d\)，故 \(\rho_d\) 良定义。它显然保持加法与乘法，并且满射。
因此
\[
\rho_{m,d}:=\rho_d\circ\bar\pi_m:\mathcal R_m\longrightarrow \ZZ/(d\ZZ)
\]
是满半环同态，而 \(\Theta_d\) 正是其核合同。于是第 \(1\) 条与第 \(2\) 条成立。

现设 \(\Theta\) 为 \(\mathcal R_m\) 上任意半环合同。经 \(\bar\pi_m\) 输运后，得到 \(\ZZ/(M_m\ZZ)\) 上的半环合同 \(\approx\)。记
\[
S:=(\ZZ/(M_m\ZZ))/\!\approx.
\]
由于 \(\ZZ/(M_m\ZZ)\) 是交换环，对任意 \([\overline x]\in S\)，类 \([-\overline x]\) 满足
\[
[\overline x]+[-\overline x]=[\overline 0],
\]
故 \(S\) 自动成为交换环，商映射因此是环同态。于是 \(\approx\) 实际上是环合同，并由某个理想 \(I\subseteq \ZZ/(M_m\ZZ)\) 唯一决定。

\(\ZZ/(M_m\ZZ)\) 的全部理想都形如 \((d)/(M_m)\)，其中 \(d\mid M_m\)，并且
\[
(\ZZ/(M_m\ZZ))/((d)/(M_m))
\cong
\ZZ/(d\ZZ).
\]
故 \(\Theta\) 必为某个 \(\Theta_d\)，唯一性亦随之成立。

最后，\(\Theta_{d_1}\le \Theta_{d_2}\) 当且仅当对应理想满足
\[
(d_1)/(M_m)\subseteq (d_2)/(M_m),
\]
而这又等价于 \(d_2\mid d_1\)。因此合同格与约数格确为反同构。证毕。
\end{proof}

\begin{theorem}[折叠同余半环的非幺端同态分类与幺自同构刚性]\label{thm:xi-fold-congruence-unital-automorphism-rigidity}
\leanverified{paper\_xi\_fold\_congruence\_unital\_automorphism\_rigidity}
定义
\[
\operatorname{End}_0(\mathcal R_m)
:=
\{\phi:\mathcal R_m\to\mathcal R_m:\ \phi\ \text{保持 }0,+,\cdot\},
\]
\[
\operatorname{Idem}(\mathcal R_m)
:=
\{e\in \mathcal R_m:\ e^2=e\}.
\]
则映射
\[
\operatorname{End}_0(\mathcal R_m)\longrightarrow \operatorname{Idem}(\mathcal R_m),
\qquad
\phi\longmapsto \phi(1)
\]
是双射，其逆由
\[
e\longmapsto \phi_e,
\qquad
\phi_e(x):=ex
\]
给出。并且在组合与乘法下，
\[
\phi_e\circ \phi_f=\phi_{ef},
\]
故 \(\operatorname{End}_0(\mathcal R_m)\) 与 \(\operatorname{Idem}(\mathcal R_m)\) 自然同构为交换幂等半格。

若
\[
M_m=\prod_{i=1}^{s}p_i^{e_i},
\qquad s=\omega(M_m),
\]
则
\[
|\operatorname{Idem}(\mathcal R_m)|
=
|\operatorname{End}_0(\mathcal R_m)|
=
2^s.
\]
特别地，
\[
\Aut_{\mathrm{unital\; semiring}}(\mathcal R_m)=\{\mathrm{id}_{\mathcal R_m}\}.
\]
\end{theorem}
\begin{proof}
经 \(\bar\pi_m\) 识别后，只需在 \(A:=\ZZ/(M_m\ZZ)\) 上证明。

任取 \(\phi\in \operatorname{End}_0(A)\)，记
\[
e:=\phi(\overline 1).
\]
由乘法保持性，
\[
e^2=\phi(\overline 1)\phi(\overline 1)=\phi(\overline 1\cdot \overline 1)=\phi(\overline 1)=e,
\]
故 \(e\) 为幂等元。对任意 \(\overline k\in A\)，有
\[
\phi(\overline k)
=
\phi(k\cdot \overline 1)
=
k\phi(\overline 1)
=
\overline k\,e.
\]
因此 \(\phi=\phi_e\)，从而 \(\phi\mapsto \phi(\overline 1)\) 为单射。

反之，若 \(e^2=e\)，则
\[
\phi_e(\overline x+\overline y)=e(\overline x+\overline y)=e\overline x+e\overline y,
\]
\[
\phi_e(\overline x\,\overline y)=e\,\overline x\,\overline y
=
e^2\,\overline x\,\overline y
=
(e\overline x)(e\overline y)
=
\phi_e(\overline x)\phi_e(\overline y),
\]
故 \(\phi_e\) 为半环端同态。这证明了双射。

对组合律，只需计算
\[
(\phi_e\circ \phi_f)(x)=e(fx)=(ef)x=\phi_{ef}(x).
\]

再由中国剩余分解
\[
A\cong \prod_{i=1}^{s}\ZZ/(p_i^{e_i}\ZZ),
\]
每个局部因子中的幂等元只能是 \(0\) 或 \(1\)，故总数恰为 \(2^s\)。于是
\[
|\operatorname{Idem}(A)|=|\operatorname{End}_0(A)|=2^s.
\]

若 \(\phi\) 还是幺半环自同构，则 \(\phi(\overline 1)=\overline 1\)，因而
\[
\phi(\overline x)=\overline x\cdot \overline 1=\overline x
\]
对全部 \(\overline x\in A\) 成立，故 \(\phi=\mathrm{id}\)。经 \(\bar\pi_m\) 传回 \(\mathcal R_m\) 即得结论。证毕。
\end{proof}

\begin{theorem}[nilradical 的显式理想闭式、基数公式与约化商]\label{thm:xi-fold-congruence-nilradical-reduced-quotient}
设
\[
M_m=\prod_{i=1}^{s}p_i^{e_i},
\qquad
\operatorname{rad}(M_m):=\prod_{i=1}^{s}p_i,
\qquad
\mathfrak N_m:=\sqrt{0}\subset \mathcal R_m.
\]
则有显式描述
\[
\mathfrak N_m
=
\{x\in \mathcal R_m:\ \bar\pi_m(x)\in \operatorname{rad}(M_m)\,\ZZ/(M_m\ZZ)\}
=
\operatorname{rad}(M_m)\,\mathcal R_m.
\]
等价地，若 \(\bar\pi_m(x)=\overline a\)，则
\[
x\in \mathfrak N_m
\iff
\operatorname{rad}(M_m)\mid a.
\]
并且：
\begin{enumerate}
  \item \(\mathfrak N_m\) 的基数精确为
  \[
  |\mathfrak N_m|=\frac{M_m}{\operatorname{rad}(M_m)};
  \]
  \item \(\mathcal R_m\) 为约化环当且仅当 \(M_m\) 平方自由；
  \item 约化商具有闭式
  \[
  \mathcal R_m/\mathfrak N_m
  \cong
  \ZZ/(\operatorname{rad}(M_m)\ZZ)
  \cong
  \prod_{i=1}^{s}\FF_{p_i}.
  \]
\end{enumerate}
\end{theorem}
\begin{proof}
经 \(\bar\pi_m\) 识别后，只需在
\[
A:=\ZZ/(M_m\ZZ)\cong \prod_{i=1}^{s}\ZZ/(p_i^{e_i}\ZZ)
\]
上证明。

先看单个局部因子 \(A_i:=\ZZ/(p_i^{e_i}\ZZ)\)。若 \(p_i\mid x\)，则 \(x^{e_i}=0\) 于 \(A_i\) 中，故 \(x\) 幂零；反之，若 \(p_i\nmid x\)，则 \(x\) 为单位，不可能幂零。故
\[
\sqrt{0}_{A_i}=p_iA_i.
\]
于是有限直积环的 nilradical 满足
\[
\sqrt{0}_{A}
\cong
\prod_{i=1}^{s}p_iA_i.
\]
经 CRT 识别，这恰是理想
\[
\operatorname{rad}(M_m)\,A.
\]
传回 \(\mathcal R_m\) 即得
\[
\mathfrak N_m=\operatorname{rad}(M_m)\,\mathcal R_m.
\]
而 \(x\in \mathfrak N_m\) 当且仅当其像 \(\bar\pi_m(x)=\overline a\) 落在该理想中，这又等价于每个 \(p_i\mid a\)，亦即 \(\operatorname{rad}(M_m)\mid a\)。

其基数为
\[
|\mathfrak N_m|
=
\prod_{i=1}^{s}|p_iA_i|
=
\prod_{i=1}^{s}p_i^{e_i-1}
=
\frac{\prod_i p_i^{e_i}}{\prod_i p_i}
=
\frac{M_m}{\operatorname{rad}(M_m)}.
\]

\(\mathcal R_m\) 约化当且仅当 \(\mathfrak N_m=0\)，这等价于每个 \(e_i=1\)，亦即 \(M_m\) 平方自由。

最后，
\[
\mathcal R_m/\mathfrak N_m
\cong
\prod_{i=1}^{s}\bigl(\ZZ/(p_i^{e_i}\ZZ)\bigr)\Big/p_i\,\ZZ/(p_i^{e_i}\ZZ)
\cong
\prod_{i=1}^{s}\FF_{p_i}.
\]
而 \(\operatorname{rad}(M_m)\) 本身平方自由，再由 CRT 得
\[
\ZZ/(\operatorname{rad}(M_m)\ZZ)\cong \prod_{i=1}^{s}\FF_{p_i}.
\]
证毕。
\end{proof}

\begin{theorem}[折叠同余半环的四相算术几何分类]\label{thm:xi-fold-congruence-crt-primitive-idempotents}
设
\[
M_m=\prod_{i=1}^{s}p_i^{e_i}
\]
为 \(M_m\) 的素因子分解，其中 \(p_i\) 两两不同。则
\[
\mathcal R_m\ \cong\ \prod_{i=1}^{s}\ZZ/(p_i^{e_i}\ZZ),
\]
并且恰有以下四种互斥且穷尽的算术几何相：
\begin{enumerate}
  \item \textbf{域相}：
  \[
  \mathcal R_m\ \text{是域}
  \iff
  M_m\ \text{是素数}.
  \]
  \item \textbf{局部非约化相}：
  \[
  \mathcal R_m\ \text{是局部环但不是域}
  \iff
  M_m=p^a,\ a\ge 2.
  \]
  \item \textbf{分裂约化相}：
  \[
  \mathcal R_m\ \text{约化且非局部}
  \iff
  s\ge 2\ \text{且}\ e_1=\cdots=e_s=1.
  \]
  在此情形下，
  \[
  \mathcal R_m\cong \prod_{i=1}^{s}\FF_{p_i}.
  \]
  \item \textbf{分裂非约化相}：
  \[
  \mathcal R_m\ \text{既非局部又非约化}
  \iff
  s\ge 2\ \text{且存在}\ i\ \text{使}\ e_i\ge 2.
  \]
\end{enumerate}
此外，\(\mathcal R_m\) 的原始中心幂等元恰有 \(s\) 个，而全部幂等元总数恰为 \(2^s\)。
\end{theorem}
\begin{proof}
Chinese remainder theorem 给出
\[
\ZZ/(M_m\ZZ)\cong \prod_{i=1}^{s}\ZZ/(p_i^{e_i}\ZZ),
\]
再与 \(\bar\pi_m\) 复合即得所示分解。

若 \(s=1\) 且 \(e_1=1\)，则
\[
\mathcal R_m\cong \ZZ/(p_1\ZZ)=\FF_{p_1},
\]
故 \(\mathcal R_m\) 为域；反之，\(\ZZ/(n\ZZ)\) 为域当且仅当 \(n\) 为素数，故得第 \(1\) 条。

若 \(s=1\) 且 \(e_1\ge 2\)，则
\[
\mathcal R_m\cong \ZZ/(p_1^{e_1}\ZZ)
\]
为局部环，其唯一极大理想是 \((p_1)/(p_1^{e_1})\)；同时 \(p_1\) 的像是非零幂零元，故它不是域。反之，若 \(\mathcal R_m\) 是局部环但不是域，则其 CRT 分解只能含一个局部因子，且该因子指数至少为 \(2\)。这证明第 \(2\) 条。

若 \(s\ge 2\) 且一切 \(e_i=1\)，则
\[
\mathcal R_m\cong \prod_{i=1}^{s}\FF_{p_i},
\]
它显然约化且非局部，从而得到第 \(3\) 条。反之，若 \(\mathcal R_m\) 约化且非局部，则必须至少有两个 CRT 因子，并且每个因子都约化，这就强制全部 \(e_i=1\)。

若 \(s\ge 2\) 且某个 \(e_i\ge 2\)，则 CRT 分解中至少含两个局部因子，所以 \(\mathcal R_m\) 非局部；同时该 \(i\)-因子含非零幂零元，故整体非约化。这证明第 \(4\) 条。反之，若 \(\mathcal R_m\) 既非局部又非约化，则 CRT 分解至少含两个因子，且至少一个因子指数 \(\ge 2\)。

上述四种情形显然互斥且穷尽。

最后，在交换环中一切幂等元都中心。每个局部因子中的幂等元只能是 \(0\) 或 \(1\)，故全体幂等元可由各坐标独立选取 \(0/1\) 得到，总数为 \(2^s\)。而原始中心幂等元正是各坐标投影
\[
e_i:=(0,\dots,0,1,0,\dots,0)
\qquad(1\le i\le s),
\]
它们两两正交且和为 \(1\)，因此恰有 \(s\) 个。证毕。
\end{proof}
\leanverified{paper\_xi\_fold\_congruence\_crt\_primitive\_idempotents}

\begin{theorem}[\texorpdfstring{$p$}{p}-primary 商塔的存在唯一性与最大可见深度]\label{thm:xi-fold-congruence-pprimary-quotient-depth}
对任意素数 \(p\)，定义
\[
\nu_p(m):=v_p(F_{m+2})=v_p(M_m).
\]
则对每个整数 \(s\ge 1\)，存在保持幺元的满半环同态
\[
\rho_{m,p^s}:\mathcal R_m\twoheadrightarrow \ZZ/(p^s\ZZ)
\]
当且仅当
\[
s\le \nu_p(m).
\]
并且一旦存在，该同态唯一，并由
\[
\rho_{m,p^s}=\mathrm{red}_{p^s}\circ \bar\pi_m
\]
给出，其中 \(\mathrm{red}_{p^s}\) 表示模 \(p^s\) 约化。换言之，\(\mathcal R_m\) 的全部 \(p\)-进可见深度完全由单个整数 \(\nu_p(m)\) 控制，而不存在额外的隐藏 \(p\)-层。
\end{theorem}
\begin{proof}
由 \(\bar\pi_m\) 的同构性，只需讨论何时存在幺半环满同态
\[
\phi:\ZZ/(M_m\ZZ)\twoheadrightarrow \ZZ/(p^s\ZZ).
\]
若 \(\phi\) 存在，则 \(\phi(\overline 1)=\overline 1\)，故对任意 \(n\in\NN\) 都有
\[
\phi(\overline n)=\overline n \pmod{p^s}.
\]
因此 \(\phi\) 若存在便被唯一确定；它良定义的充要条件是
\[
\overline{M_m}\mapsto \overline{0}\ \text{于}\ \ZZ/(p^s\ZZ),
\]
即
\[
p^s\mid M_m.
\]
这等价于 \(s\le \nu_p(m)\)。

若 \(p^s\mid M_m\)，则自然约化
\[
\mathrm{red}_{p^s}:\ZZ/(M_m\ZZ)\longrightarrow \ZZ/(p^s\ZZ)
\]
良定义且满射，复合 \(\bar\pi_m\) 即得所需的 \(\rho_{m,p^s}\)。唯一性已由 \(\phi(\overline 1)=\overline 1\) 的强制性给出。证毕。
\end{proof}

\begin{corollary}[规范二值商的存在判据与 dyadic 深度]\label{cor:xi-fold-congruence-binary-quotient-period3}
\leanverified{paper\_xi\_fold\_congruence\_binary\_quotient\_period3}
存在唯一的保持幺元满半环同态
\[
\mathcal R_m\twoheadrightarrow \ZZ/(2\ZZ)
\]
当且仅当
\[
3\mid (m+2).
\]
并且 \(\mathcal R_m\) 的最大 dyadic 可见深度恰为
\[
\nu_2(m)=v_2(F_{m+2}).
\]
\end{corollary}
\begin{proof}
由定理 \ref{thm:xi-fold-congruence-pprimary-quotient-depth} 取 \(p=2\)、\(s=1\)，上述二值商存在当且仅当 \(2\mid F_{m+2}\)。
而 Fibonacci 数列模 \(2\) 满足周期 \(3\)：
\[
F_1\equiv 1,\quad F_2\equiv 1,\quad F_3\equiv 0,\quad
F_{n+3}\equiv F_n\pmod 2.
\]
故
\[
2\mid F_n\iff 3\mid n.
\]
取 \(n=m+2\) 即得第一条。第二条则是定理
\ref{thm:xi-fold-congruence-pprimary-quotient-depth} 在 \(p=2\) 时的直接专门化。证毕。
\end{proof}

由此，折叠同余层上的全部半环合同、全部非幺端同态、nilradical 厚度、四相算术几何与全部 \(p\)-primary 商塔都退化为单个 Fibonacci 模数 \(M_m\) 的约数结构、幂等结构与素因子型。

\endinput
